\section{Introduction}

Multiplication is usually treated as scalar arithmetic: two integers enter and
one integer exits. In a positional base, however, multiplication has internal
geometry. Digits occupy degree positions, digit pairs interact across a
two-dimensional index grid, and carry propagation moves excess mass along a
one-dimensional degree axis. Projection--Carry Geometry makes this structure
explicit.

Fix a base \(B\ge 2\). We write nonnegative integers as digit polynomials
\[
    X=\sum_i a_i B^i,\qquad
    Y=\sum_j b_j B^j,
\]
where \(0\le a_i,b_j<B\). The product is decomposed into three stages. First,
the digit interaction field is the outer-product matrix
\[
    M_{ij}=a_i b_j.
\]
Second, anti-diagonal projection aggregates interactions with the same total
degree,
\[
    c_k=\sum_{i+j=k}M_{ij}.
\]
Third, the raw coefficients \(c_k\) are converted into valid base-\(B\) digits
by carry normalization.

The main thesis is that base-\(B\) multiplication decomposes into bilinear
diagonal projection followed by nonlinear carry flow. The projection step is
ordinary convolution of digit polynomials. The carry step is a directed mass
flow along the degree axis, locally activated when \(c_k+r_k\ge B\), where
\(r_k\) is incoming carry.

This viewpoint separates potential interaction geometry from realized carry
behavior. If \(S_X=\{i:a_i\ne 0\}\) and \(S_Y=\{j:b_j\ne 0\}\), then the raw
support before carry is the sumset \(S_X+S_Y\). The representation counts
\[
    \rho_{X,Y}(k)=\#\{(i,j)\in S_X\times S_Y:i+j=k\}
\]
measure concentration of possible interactions on anti-diagonals. We argue
that this support-sumset geometry shapes carry complexity, while digit
amplitudes and incoming carry determine the realized local events.

The contributions are:
\begin{itemize}
    \item[C1.] A matrix projection formulation of positional multiplication.
    \item[C2.] A formulation of carry as nonlinear mass flow.
    \item[C3.] A local carry activation law.
    \item[C4.] Support-sumset geometry and representation-count bounds for raw
    diagonal mass.
    \item[C5.] Carry-complexity metrics and empirical scaling experiments.
\end{itemize}

This is a structural and geometric framework, not a new fast multiplication
algorithm. Its purpose is to expose the anatomy of multiplication and carry
propagation, not to replace existing arithmetic algorithms.
